% Generated draft: requires mathematical and source review.
\documentclass[11pt]{article}
\usepackage[margin=25mm]{geometry}
\usepackage{fontspec}
\setmainfont{Latin Modern Roman}
\usepackage{amsmath,amssymb,mathtools,unicode-math}
\setmathfont{Latin Modern Math}
\usepackage{enumitem,needspace,xurl,hyperref}
\setlength{\parindent}{0pt}
\setlength{\parskip}{5pt}
\setlength{\emergencystretch}{3em}
\newcommand{\N}{\mathbb{N}}
\newcommand{\Z}{\mathbb{Z}}
\newcommand{\Q}{\mathbb{Q}}
\newcommand{\R}{\mathbb{R}}
\newcommand{\C}{\mathbb{C}}
\newcommand{\F}{\mathbb{F}}
\newcommand{\A}{\mathbb{A}}
\newcommand{\PP}{\mathbb{P}}
\newcommand{\E}{\mathbb{E}}
\newcommand{\eps}{\varepsilon}
\newcommand{\dd}{\,\mathrm{d}}
\begin{document}
\section*{UnsolvedMath ID: 700730}
\textbf{Source problem code:} catalogue entry
\subsection{Definitions and mathematical statement}
For positive integers \(m,n\), let \(F_m\) be the free group on generators \(x_1,\ldots,x_m\), and define
\[
N_n=\left\langle\!\left\langle w^n\mid w\in F_m\right\rangle\!\right\rangle_{F_m},
\qquad B(m,n)=F_m/N_n.
\]
Here \(\langle\!\langle S\rangle\!\rangle_{F_m}\) denotes the normal closure of a subset \(S\subseteq F_m\). Write \(\overline v\) for the image of \(v\in F_m\) in \(B(m,n)\). Positive finite rank and positive integer exponent are the conventions used here; the supplied source does not explicitly specify the ranges of \(m\) and \(n\).

Define the presentation property \(\mathcal P(m,n)\) to mean that there exists a set \(V\subseteq F_m\) such that
\[
N_n=\left\langle\!\left\langle v^n\mid v\in V\right\rangle\!\right\rangle_{F_m}
\]
and
\[
\forall v\in V\;\forall d\in\mathbb Z_{>0},\qquad
\bigl(d\mid n\ \text{and}\ d<n\bigr)
\Longrightarrow \overline v^{\,d}\ne 1.
\]
Thus the required defining relations present \(B(m,n)\) on its standard free generators, and the base word of each defining relation has order exactly \(n\) in \(B(m,n)\). No finiteness requirement on \(V\) is imposed by the source.

\begin{enumerate}
\item[(a)] Does \(\mathcal P(m,n)\) hold for every positive integer rank \(m\) and every positive integer exponent \(n\)? The supplied statement reports that it holds for odd \(n\geq665\), and also for all \(n\geq2^{48}\) satisfying \(2^9\mid n\), with arbitrary \(m\) as in the question. These are source-reported qualifications, not independently verified assertions here.
\item[(b)] Does the same property \(\mathcal P(m,n)\) hold for arbitrary \(m\) when \(n=2^\ell\gg1\)? Here \(\ell\) is a nonnegative integer, and the source notation means sufficiently large powers of two. The source supplies neither a threshold nor an explicit statement of whether that threshold must be independent of \(m\); this quantifier qualification remains unresolved.
\end{enumerate}

The supplied record places part (b) in its archive of solved parts and reports an affirmative answer attributed to S. V. Ivanov. It labels the overall problem \texttt{partially\_solved}, with remaining part(s) in the main section. These are supplied archival and status assessments, not independent mathematical verification.
\subsection{Short English statement}
Can every free Burnside group \(B(m,n)\) be presented using only relations \(v^n=1\) whose base words represent elements of exact order \(n\)? The source reports two established exponent ranges and an affirmative archived answer for sufficiently large powers of two. The meaning of sufficiently large, including its possible dependence on \(m\), is unspecified.
\subsection{Sources}
\begin{itemize}
\item Local UnsolvedMath source registry, permanent ID \(700730\), exact source problem code \texttt{catalogue entry}, title \emph{Kourovka Notebook Problem 11.37}. The supplied upstream revision is \texttt{d371ee917ac9647fbb06cd6fc529deec6f3dda82}.
\item \emph{The Kourovka Notebook (Unsolved Problems in Group Theory)}, problem 11.37, proposed by S. V. Ivanov in the 11th issue (1990); the supplied background identifies an update in the 21st issue (2026), arXiv:1401.0300v48. The supplied collection description names E. I. Khukhro and V. D. Mazurov as editors. These references have not been independently consulted here.
\item Supplied Archive of Solved Problems attribution for the affirmative answer to part (b): S. V. Ivanov, \emph{Int. J. Algebra Comput.}, volume 4, nos. 1--2 (1994), pp. 1--308. No article title or direct link is supplied.
\item Supplied status review checked on \texttt{2026-10-16}: canonical status \texttt{partially\_solved}, evidence field \texttt{published}; its note assigns solved part(s) to the archive and remaining part(s) to the main section. This is supplied dated curation, not an independently verified literature assessment.
\end{itemize}
\end{document}
